\title{

\begin{singlespacing}

\large

\vspace{-100pt}

\rule{\textwidth}{2pt}
\vfill\vfill
{\Huge\textbf{\thesistitle}}
\vfill
\rule{\textwidth}{2pt}

\vfill

by

\vfill

{\LARGE\authorname}

\vfill

\includegraphics[width=0.5\textwidth]{figures/frontmatter/logo}
\\
A thesis submitted to the {\university} for the degree of
\\
\textsc{\degreetitle}

\vfill

\begin{flushright}
\groupname
\\
\department
\\
\faculty
\\
\university
\vfill
\monthyear\today
\end{flushright}

\date{}

\end{singlespacing}

}
